\documentclass[10pt,a4paper]{amsart}
\usepackage[margin=19mm]{geometry}
\usepackage{amsmath,amssymb,mathtools}
\usepackage[scr=rsfs,cal=euler]{mathalfa}
\usepackage{xfrac}
\usepackage[unicode,hidelinks,bookmarks=false]{hyperref}
\newcommand{\ie}{i.e.\ }
\newcommand{\cf}{cf.\ }
\newcommand{\resp}{respectively\ }
\newcommand{\Subsection}{\subsection}
\providecommand{\nobreakdash}{\nobreak\hbox{-}}
\setlength{\emergencystretch}{2em}
\multlinegap0pt
\usepackage{bm}
\usepackage{bbm}
\usepackage{dsfont}
\usepackage{xspace}
\usepackage{cancel}
\usepackage{mathtools}
\usepackage{amsthm}
\makeatletter
\@ifundefined{theorem}{\newtheorem{theorem}{Theorem}}{}
\@ifundefined{proposition}{\newtheorem{proposition}[theorem]{Proposition}}{}
\makeatother
\title[Zero blocking numbers of graphs with complexity results]{Zero blocking numbers of graphs with complexity results}
\author{Hau-Yi Lin, Wu-Hsiung Lin, Gerard Jennhwa Chang}
\date{Source: arXiv:2508.17785v1 (CC BY 4.0)}
\begin{document}
\maketitle

\noindent\emph{Source and licence.} Zero blocking numbers of graphs with complexity results. Version arXiv:2508.17785v1 (CC BY 4.0), distributed under the Creative Commons Attribution 4.0 International licence, as recorded in the arXiv licence metadata for this exact version and shown on the abstract page https://arxiv.org/abs/2508.17785v1. Full rights notice: licenses/arxiv-2508.17785-CC-BY-4.0.txt.\par\smallskip

\noindent\textit{Editorial scope.} Counted unit: the Theorem whose source label is \texttt{join} in the article (source0.tex lines 364--375, source label \texttt{join}), delivered together with the article's own complete original proof of it (source0.tex lines 376--448, one \texttt{proof} environment of 73 lines). No mathematics is written, repaired or re-derived: statement and proof are the article's own text line for line. Required context retained uncounted: 2nd (lines 341--351). Every definition, display and label of those lines is retained unchanged. Omitted from this package and not redistributed: 1--340, 352--363, 449--883. Nothing inside a retained excerpt is omitted or altered: every retained body is delivered exactly as the article's lines are written, so no omission receipt of any kind accompanies this package. Citations: the retained excerpts carry no citation command at all, so no bibliography entry of the article is transcribed and no reference is redistributed. Reference adaptation: none was needed and none was made; every reference inside the delivered unit resolves inside it, so no unresolved reference or citation is produced. Printed slips kept literal: no printed word, symbol, hypothesis or number of the counted statement or of its proof was altered. Layout and class: the article's own class is \texttt{elsarticle} with packages including amssymb, amsmath, amsthm, xcolor, graphicx, tikz, algorithm, algorithmicx, algpseudocode; the wrapper is the lane's pinned \texttt{amsart} editorial wrapper at 10pt on A4, which restates the theorem-like environments the retained text needs and adds the article's own macro definitions that the retained text needs (none), with extra wrapper packages (bm, bbm, dsfont, xspace, cancel, mathtools). The delivered numbering is the unit's own. Assets: no retained span contains a figure environment, a graphics-inclusion command, a PDF-page inclusion, a table or a TikZ picture; no image file is copied into this package, the package declares no asset dependency and no image file is redistributed with it. Licence: version arXiv:2508.17785v1 of this article is distributed under the Creative Commons Attribution 4.0 International licence, as recorded in the arXiv licence metadata for this exact version and shown on the abstract page https://arxiv.org/abs/2508.17785v1. A full rights notice is delivered at licenses/arxiv-2508.17785-CC-BY-4.0.txt. Characters: every retained excerpt is pure ASCII, so the delivered engine, XeLaTeX with its default Latin Modern text font, reports no missing character.
\begin{proposition}\label{2nd}
For a graph $G$, we have 
$$
B'(G)=\left\{\begin{array}{ll}
         2,      & \mbox{if $G$ has at least two isolated vertices}; \\
         B(G-x), & \mbox{if $G$ has exactly one isolated vertex $x$}; \\
         B(G),   & \mbox{if $G$ has no isolated vertex}.
             \end{array}
       \right.
$$ 
\end{proposition}

\begin{theorem} \label{join}
If $G$ and $H$ are graphs of $n$ and $m$ vertices respectively, then
$B'(G+H)=B(G+H)=\min\{B'(G),B'(H),a\}$, where
$$
a=\left\{\begin{array}{ll}
 \gamma(G)+\gamma(H), & \mbox{if $n=1$ or $m=1$ or $\gamma(G)=\gamma(H)=1$}; \\
 3,                   & \mbox{if $n\ge 2$, $m\ge 2$ and}\\ 
                      & \mbox{$(\gamma(G)\le 2$ or $\gamma(H)\le 2$ but not $\gamma(G)=\gamma(H)=1)$}; \\
 4,                   & \mbox{if $\gamma(G)\ge 3$ and $\gamma(H)\ge 3$}.
  \end{array}\right.
$$
\end{theorem}

\begin{proof}
Since $G+H$ has no isolated vertex, $B'(G+H)=B(G+H)$ by Proposition \ref{2nd}.

First, prove that $B(G+H)\le \min\{B'(G),B'(H),a\}$.
Since a minimum second zero blocking set of $G$ is
a zero blocking set of $G+H$, we have $B(G+H)\le B'(G)$.
Similarly, $B(G+H)\le B'(H)$.
Now we construct a zero blocking set of $G+H$ for each condition of $a$ in three cases.

Case 1. $n=1$ or $m=1$ or $\gamma(G)=\gamma(H)=1$.
In this case we choose
a minimum dominating set $D$ (respectively, $D'$) of $G$ (respectively, $H$).
Then $D\cup D'$ is a zero blocking set of $G+H$,
and so $B(G+H)\le |D\cup D'|= \gamma(G)+\gamma(H)=a$, 
the last equality follows from the definition if $a$.

Case 2. $n\ge 2$, $m\ge 2$,
($\gamma(G)\le 2$ or $\gamma(H)\le 2$ but not $\gamma(G)=\gamma(H)=1$).
In this case, $a=3$.
By symmetry, assume $\gamma(G)\le 2$.
As $n\ge 2$, we may choose
a dominating set $D$ of $G$ with $|D|=2$ and a vertex $y\in V(H)$.
Then $D\cup\{y\}$ is a zero blocking set of $G+H$,
and so $B(G+H)\le |D\cup\{y\}| = 3=a$.

Case 3. $\gamma(G)\ge 3$ and $\gamma(H)\ge 3$.
In this case, $a=4$. As $n\ge 3$ and $m\ge 3$, we may choose vertices $x,y\in V(G)$ and $z,w\in V(H)$.
Then $\{x,y,z,w\}$ is a zero blocking set of $G+H$, and so $B(G+H)\le |\{x,y,z,w\}|=4=a$.

In summary, $B(G+H)\le \min\{B'(G),B'(H),a\}$.

To show $B(G+H)\ge \min\{B'(G),B'(H),a\}$,
we choose a minimum zero blocking set $W$ of $G+H$.
Let $D=W\cap V(G)$ and $D'=W\cap V(H)$.
If $|D'|=0$, then every vertex in $V(H)$ is black and so has at least two white neighbors, which are in $V(G)$.
Also, $W=D$ is a zero blocking set of $G$ and so is a second zero blocking set of $G$.
Hence $ B(G+H) =  |W|\ge B'(G)$.
Similarly, if $|D|=0$, then $B(G+H) =  |W|\ge B'(H)$.

Now suppose $|D|\ge 1$ and $|D'|\ge 1$.
We consider the following cases. First, claim that if $|D|=1$ then $D'$ is a dominating set of $H$,
for otherwise $V(H)$ has a black vertex adjacent to the only white vertex in $D$,
a contradiction.
Similarly, if $|D'|=1$, then $D$ is dominating set of $G$.

Case 1.
If $n=1$, then $|D|=1$ and so $D$ (respectively, $D'$) is a dominating set of $G$ (respectively, $H$ by the above claim).
In this case, $ B(G+H) =  |W|=|D|+|D'|\ge \gamma(G)+\gamma(H)=a$.
Similarly, if $m=1$, then $ B(G+H) = |W|\ge a$.

Case 2.
%Now suppose
If $n\geq 2$ and $m\ge 2$,
then $a=2,3$ or $4$ depending on the values of $\gamma(G)$ and $\gamma(H)$.

Case 2.1.
If $|W|=2$, then $|D|=|D'|=1$ and so $D$ (respectively, $D'$) is a dominating set of $G$ (respectively, $H$) by the above claim.
In this case, $\gamma(G)=\gamma(H)=1$, implying $a=\gamma(G)+\gamma(H)=2$.
Thus $ B(G+H) = |W|=a$.

Case 2.2.
If $|W|=3$, then either $|D|=1$ or $|D'|=1$.
By symmetry, assume that $|D|=1$ and so $|D'|=2$.
By the above claim, $D'$ is a dominating set of $H$
and so $\gamma(H)\le 2$, implying $a=2$ or $3$.
Thus $B(G+H) = |W|\ge a$.

Case 2.3.
Otherwise, $B(G+H) = |W|\ge 4\ge a$.
 
In summary, $B(G+H)\ge \min\{B'(G),B'(H),a\}$,
and theorem follows.
\end{proof}

\end{document}
